\documentclass[11pt, a4paper]{article}

% Gemeinsame Style-Definitionen (TEXINPUTS muss templates/ enthalten — siehe scripts/build_tex.py)
% bfw-style.tex — gemeinsame Style-Definitionen für alle Handouts/Aufgabenhefte
% Voraussetzung: Kompilation mit xelatex (wegen fontspec)

\usepackage[ngerman]{babel}
% Babel-ngerman macht " zu einem aktiven Shorthand (z.B. für "a -> ä). In unseren
% Texten verwenden wir " als normales Zeichen — daher Shorthand komplett ausschalten.
\shorthandoff{"}
\usepackage{geometry}
\geometry{a4paper, top=2.4cm, bottom=2.4cm, left=2.3cm, right=2.3cm,
          headheight=15pt, headsep=0.7cm, footskip=1.1cm}

\usepackage{fontspec}
% Fallback-Kette: Source Sans Pro -> Calibri -> Segoe UI -> Latin Modern Sans
% (Latin Modern ist immer mit MiKTeX vorhanden)
\IfFontExistsTF{Source Sans Pro}{%
  \setmainfont{Source Sans Pro}[Ligatures=TeX]%
}{\IfFontExistsTF{Calibri}{%
  \setmainfont{Calibri}[Ligatures=TeX]%
}{\IfFontExistsTF{Segoe UI}{%
  \setmainfont{Segoe UI}[Ligatures=TeX]%
}{%
  \setmainfont{Latin Modern Sans}[Ligatures=TeX]%
}}}
\IfFontExistsTF{Source Code Pro}{%
  \setmonofont{Source Code Pro}[Scale=0.92]%
}{\IfFontExistsTF{Consolas}{%
  \setmonofont{Consolas}[Scale=0.92]%
}{%
  \setmonofont{Latin Modern Mono}[Scale=0.92]%
}}

% Gesperrte Versalien für Labels/Titelbalken (Kapitälchen-Ersatz, funktioniert
% mit jeder OpenType-Schrift — Latin Modern Sans hat keine echten Kapitälchen)
\newcommand{\bfwlabelfont}{\footnotesize\bfseries\addfontfeatures{LetterSpace=8.0}}

\usepackage{xcolor}
% Farbpalette: hell, freundlich, modern
\definecolor{bfwPrimary}{HTML}{0EA5A4}    % Petrol/Türkis - Primär
\definecolor{bfwPrimaryDark}{HTML}{0F766E} % dunkler Petrol für Headings
\definecolor{bfwAccent}{HTML}{F59E0B}     % Amber - Akzent
\definecolor{bfwSuccess}{HTML}{10B981}    % Grün - Erfolg / Lösung
\definecolor{bfwWarn}{HTML}{F97316}       % Orange - Achtung
\definecolor{bfwInk}{HTML}{0F172A}        % fast Schwarz
\definecolor{bfwInkSoft}{HTML}{475569}    % gedämpftes Grau
\definecolor{bfwBgSoft}{HTML}{F8FAFC}     % off-white
\definecolor{bfwBgTeal}{HTML}{ECFEFF}     % helles Türkis
\definecolor{bfwBgAmber}{HTML}{FEF3C7}    % helles Gelb
\definecolor{bfwBgRose}{HTML}{FFE4E6}     % helles Rosé für Eselsbrücken
\definecolor{bfwTabHead}{HTML}{E4F3F2}    % zarte Kopfzeilen-Hinterlegung für Tabellen

\usepackage[colorlinks=true, linkcolor=bfwPrimaryDark, urlcolor=bfwPrimaryDark]{hyperref}
\usepackage{lastpage}
\usepackage{enumitem}
\setlist[itemize]{leftmargin=*, itemsep=2pt, topsep=2pt}
\setlist[enumerate]{leftmargin=*, itemsep=2pt, topsep=2pt}

\usepackage[most]{tcolorbox}
\usepackage{booktabs}
\usepackage{colortbl}
\usepackage{tikz}
\usetikzlibrary{decorations.pathmorphing, positioning, shapes.geometric, arrows.meta}

% ---------------------------------------------------------------------------
% Einheitliches Tabellen-Design
%   - etwas mehr Zeilenluft, dezente graue Linien (wirkt auch mit \hline ruhiger)
%   - Kopfzeile einer Tabelle bei Bedarf zart hinterlegen: \rowcolor{bfwTabHead}
% ---------------------------------------------------------------------------
\renewcommand{\arraystretch}{1.25}
\arrayrulecolor{bfwInkSoft!50}
\setlength{\heavyrulewidth}{1pt}
\setlength{\lightrulewidth}{0.5pt}

% ---------------------------------------------------------------------------
% Boxen — klare farbige Titelbalken mit gesperrten Versal-Labels
% (Titel bleibt per Option überschreibbar: \begin{merksatz}[title={...}])
% ---------------------------------------------------------------------------
\tcbset{bfwbox/.style={%
  enhanced, breakable,
  boxrule=0.5pt, arc=2.5pt,
  left=10pt, right=10pt, top=8pt, bottom=8pt,
  toptitle=4pt, bottomtitle=4pt,
  titlerule=0pt,
  fonttitle=\bfwlabelfont,
  coltitle=white
}}

% Merkkästchen — wichtige Definition / Kernpunkt
\newtcolorbox{merksatz}[1][]{%
  bfwbox,
  colback=bfwBgTeal,
  colframe=bfwPrimaryDark,
  colbacktitle=bfwPrimaryDark,
  title={MERKE},
  #1
}

% Eselsbrücke — Mnemoniks
\newtcolorbox{eselsbruecke}[1][]{%
  bfwbox,
  colback=bfwBgRose,
  colframe=bfwAccent!85!black,
  colbacktitle=bfwAccent!85!black,
  title={ESELSBRÜCKE},
  #1
}

% Beispiel-Box
\newtcolorbox{beispiel}[1][]{%
  bfwbox,
  colback=bfwBgSoft,
  colframe=bfwInkSoft,
  colbacktitle=bfwInkSoft,
  title={BEISPIEL},
  #1
}

% Lösungs-Box (für Aufgabenhefte)
\newtcolorbox{loesung}[1][]{%
  bfwbox,
  colback=bfwSuccess!7,
  colframe=bfwSuccess!65!black,
  colbacktitle=bfwSuccess!65!black,
  title={LÖSUNG},
  #1
}

% Achtung-Box — typische Fehler / Stolperfallen
\newtcolorbox{achtung}[1][]{%
  bfwbox,
  colback=bfwWarn!6,
  colframe=bfwWarn!85!black,
  colbacktitle=bfwWarn!85!black,
  title={ACHTUNG},
  #1
}

% Formel-Box — für zentrale Formeln (groß, ruhig, ohne Titelbalken)
\newtcolorbox{formelbox}[1][]{%
  enhanced,
  colback=bfwBgTeal,
  colframe=bfwPrimaryDark,
  boxrule=1pt, arc=3pt,
  left=10pt, right=10pt, top=6pt, bottom=6pt,
  #1
}

% Glossar-Eintrag
\newcommand{\glossareintrag}[2]{%
  \par\noindent\textbf{\color{bfwPrimaryDark}#1} \enspace #2\par\smallskip
}

% ---------------------------------------------------------------------------
% Abschnitts-Design: farbiges Nummern-Quadrat + feine Linie unter \section,
% dezent farbige \subsection/\subsubsection
% ---------------------------------------------------------------------------
\usepackage{titlesec}
\newcommand{\bfwSecBadge}[1]{%
  \begin{tikzpicture}[baseline={([yshift=-0.62em]n.center)}]
    \node[fill=bfwPrimaryDark, text=white, font=\normalsize\bfseries,
          minimum width=1.55em, minimum height=1.55em, inner sep=1pt,
          rounded corners=1.5pt] (n) {#1};
  \end{tikzpicture}}
\titleformat{\section}
  {\Large\bfseries\color{bfwPrimaryDark}}
  {\bfwSecBadge{\thesection}}{0.6em}{}
  [\vspace{-0.2em}{\color{bfwPrimary!60}\titlerule[0.8pt]}]
\titleformat{\subsection}
  {\large\bfseries\color{bfwPrimaryDark}}
  {\textcolor{bfwPrimary}{\thesubsection}}{0.5em}{}
\titleformat{\subsubsection}
  {\normalsize\bfseries\color{bfwInk}}
  {\textcolor{bfwPrimary}{\thesubsubsection}}{0.4em}{}
\titlespacing*{\section}{0pt}{1.6em}{1em}
\titlespacing*{\subsection}{0pt}{1.2em}{0.5em}

% TikZ-Stil "skizzenhaft" — etwas weicher als Rough.js, aber stabil
\tikzset{
  roughbox/.style = {
    draw=bfwPrimaryDark, thick, fill=bfwBgTeal,
    rounded corners=4pt, inner sep=6pt
  },
  roughaccent/.style = {
    draw=bfwAccent, thick, fill=bfwBgAmber,
    rounded corners=4pt, inner sep=6pt
  },
  roughline/.style = {
    draw=bfwInkSoft, thick, -{Stealth[length=2.5mm]}
  }
}

% Bullet-Symbole (bewusst kein FontAwesome — vermeidet Font-Installations-Probleme)
\newcommand{\faLightbulb}{$\bullet$}
\newcommand{\faBookmark}{$\bullet$}

% ---------------------------------------------------------------------------
% Kopf-/Fußzeile
%   Kopf links:  Wortlogo „Wirtschaftsfachwirt"
%   Kopf rechts: Dokument-Kurztext, gesetzt via \kopfzeile{...}
%   Fuß links:   © Can Siebert · 2026 — Fuß rechts: Seite X von Y
%   Titelseite (Seite 1) bleibt ohne Kopf-/Fußzeile (\handouttitel setzt
%   \thispagestyle{empty}).
% ---------------------------------------------------------------------------
\usepackage{fancyhdr}
\newcommand{\bfwKopfzeileText}{}
\newcommand{\kopfzeile}[1]{\renewcommand{\bfwKopfzeileText}{#1}}
\pagestyle{fancy}
\fancyhf{}
\fancyhead[L]{\small\bfseries\color{bfwPrimaryDark}Wirtschaftsfachwirt}
\fancyhead[R]{\small\color{bfwInkSoft}\bfwKopfzeileText}
\renewcommand{\headrulewidth}{0.6pt}
\renewcommand{\headrule}{{\color{bfwPrimary}\hrule width\headwidth height 0.6pt}}
\fancyfoot[L]{\small\color{bfwInkSoft}\copyright\ Can Siebert\,\textperiodcentered\,2026}
\fancyfoot[R]{\small\color{bfwInkSoft}Seite \thepage\ von \pageref*{LastPage}}
% Auch \thispagestyle{plain} (z.B. durch \maketitle) soll leer bleiben:
\fancypagestyle{plain}{\fancyhf{}\renewcommand{\headrulewidth}{0pt}\renewcommand{\headrule}{}}

% ---------------------------------------------------------------------------
% \handouttitel{Obertitel}{Untertitel}{Kurs-Label}{Termin-Zeile}
%   Einheitlicher professioneller Titelblock: Dachzeile, farbige Headline,
%   Untertitel, farbige Trennlinie und dezente Meta-Zeile (Kurs/Termin/Dozent).
%   Ersetzt \title/\maketitle. Direkt nach \begin{document} aufrufen.
% ---------------------------------------------------------------------------
\newcommand{\handouttitel}[4]{%
  \thispagestyle{empty}%
  \begingroup
  \setlength{\parindent}{0pt}%
  {\bfwlabelfont\color{bfwPrimary}WIRTSCHAFTSFACHWIRT (IHK)\par}
  \vspace{0.55em}
  {\huge\bfseries\color{bfwPrimaryDark}#1\par}
  \vspace{0.5em}
  {\large\color{bfwInkSoft}#2\par}
  \vspace{0.9em}
  {\color{bfwPrimary}\rule{\linewidth}{2pt}}\par
  \vspace{0.55em}
  {\small\color{bfwInkSoft}#3\ \,\textperiodcentered\,\ #4\ \,\textperiodcentered\,\ Dozent: Can Siebert\par}
  \endgroup
  \vspace{1.4em}%
}


\kopfzeile{Handout BWL Lektion 4 \textperiodcentered{} Teil 1}

\begin{document}
\handouttitel{Lektion~4 \textperiodcentered{} Teil~1: Kooperation \& Kartellrecht}%
  {Warum sich Unternehmen verbinden --- und wie sie dabei selbstständig bleiben}%
  {1.4 Unternehmenszusammenschlüsse}%
  {Sa, 27.02.2027 \textperiodcentered{} Session 1}

\begin{merksatz}[title={\bfseries Worum geht es in Teil 1?}]
Unternehmen wachsen nicht nur aus eigener Kraft (internes Wachstum), sondern auch durch
Verbindungen mit anderen Unternehmen (externes Wachstum). In dieser ersten Session
reaktivieren wir zunächst die Rechtsformen aus Lektion~3, klären dann, \emph{warum} sich
Unternehmen zusammenschließen, und grenzen \textbf{Kooperation} (Selbstständigkeit bleibt
erhalten) von \textbf{Konzentration} (wirtschaftliche Selbstständigkeit geht verloren)
ab. Anschließend lernen Sie alle \textbf{Kooperationsformen} kennen --- von der
Arbeitsgemeinschaft bis zum grundsätzlich verbotenen \textbf{Kartell} (\S\,1 GWB).
Teil~2 (Session 2) behandelt die Konzentrationsformen, die Richtungen und die
Fusionskontrolle.
\end{merksatz}

\begin{beispiel}[title={Hinweis zur Literatur}]
Der Textband-Auszug zu diesem Thema wird nachgereicht. Bis dahin dienen die beiden
Teil-Handouts als Skriptersatz; ergänzend: Übungsband Betriebswirtschaft,
Übungen 54--58 (mit Lösungen).
\end{beispiel}

\tableofcontents
\vspace{1em}
\hrule
\vspace{1em}

\section{Anknüpfung: Rechtsformen als Bausteine (Wiederholung Lektion 3)}

Zusammenschlüsse verbinden \emph{bestehende} Unternehmen --- also die Rechtsformen aus
Lektion~3. Zur Erinnerung:

\begin{itemize}
  \item \textbf{Personengesellschaften} (GbR, OHG, KG): keine juristischen Personen,
    persönliche Haftung mindestens eines Gesellschafters. Die \emph{GbR} ist die typische
    Rechtsform für Arbeitsgemeinschaften und Konsortien.
  \item \textbf{Kapitalgesellschaften} (GmbH, AG): juristische Personen, Haftung auf das
    Gesellschaftsvermögen beschränkt. Sie sind die typischen Bausteine von
    \emph{Konzernen} (Mutter- und Tochtergesellschaften) und von \emph{Joint Ventures}
    (Gemeinschafts-GmbH).
  \item Die \textbf{AG} kann Anteile anderer Gesellschaften erwerben ---
    Mehrheitsbeteiligungen sind der häufigste Weg in den Konzern.
\end{itemize}

\begin{merksatz}
Ein Zusammenschluss ändert nicht automatisch die Rechtsform: Im \emph{Konzern} bleiben
alle Gesellschaften rechtlich bestehen; erst bei der \emph{Fusion (Verschmelzung)}
erlischt mindestens ein Rechtsträger (dazu ausführlich Teil~2).
\end{merksatz}

\section{Gründe und Ziele von Unternehmenszusammenschlüssen}

Unternehmen schließen sich zusammen, um gemeinsam Ziele zu erreichen, die allein nicht
oder nur schwer erreichbar wären:

\begin{itemize}
  \item \textbf{Rationalisierung und Kostensenkung:} Größenvorteile (Kostendegression),
    gemeinsamer Einkauf (Mengenrabatte), gemeinsame Verwaltung, Vertrieb, Logistik;
  \item \textbf{Synergieeffekte:} Zusammenwirken schafft Mehrwert („2 + 2 = 5``) ---
    z.\,B. Teilung der Forschungs- und Entwicklungskosten, Wissens- und
    Technologietransfer;
  \item \textbf{Marktmacht und Wettbewerbsposition:} größere Marktanteile, Ausschalten
    von Konkurrenz, stärkere Position gegenüber Lieferanten und Abnehmern;
  \item \textbf{Sicherung von Beschaffung und Absatz:} Zugang zu Rohstoffen,
    Vorprodukten und Vertriebswegen (vertikale Motive);
  \item \textbf{Risikostreuung:} Verteilung des unternehmerischen Risikos auf mehrere
    Geschäftsfelder oder Projekte (Diversifikation);
  \item \textbf{Kapitalkraft und Finanzierung:} bessere Kreditkonditionen, gemeinsame
    Finanzierung von Großprojekten;
  \item \textbf{Schutz vor Übernahmen} und Erschließung neuer (Auslands-)Märkte.
\end{itemize}

\begin{eselsbruecke}
\textbf{SMaRRt} zusammenschließen: \textbf{S}ynergien, \textbf{Ma}rktmacht,
\textbf{R}ationalisierung, \textbf{R}isikostreuung --- die vier Kernziele, die die IHK
am liebsten abfragt.
\end{eselsbruecke}

\section{Das Schlüsselkriterium: Selbstständigkeit}

Alle Zusammenschlussformen werden danach eingeteilt, was mit der
\textbf{rechtlichen} und der \textbf{wirtschaftlichen} Selbstständigkeit der beteiligten
Unternehmen geschieht:

\begin{itemize}
  \item \textbf{Rechtliche Selbstständigkeit:} Das Unternehmen bleibt eigener
    Rechtsträger (eigene Firma, eigene Organe, eigene Bilanz).
  \item \textbf{Wirtschaftliche Selbstständigkeit:} Das Unternehmen trifft seine
    unternehmerischen Entscheidungen (Preise, Programme, Investitionen) selbst.
\end{itemize}

\subsection{Kooperation vs. Konzentration im Überblick}

\begin{center}
\begin{tabular}{p{4.2cm} p{4.9cm} p{4.9cm}}
\toprule
\textbf{Kriterium} & \textbf{Kooperation} & \textbf{Konzentration}\\
\midrule
Rechtliche Selbstständigkeit & bleibt erhalten & Konzern: bleibt erhalten;\newline Fusion/Trust: geht (bei mind. einem Partner) verloren\\
\addlinespace[3pt]
Wirtschaftliche Selbstständigkeit & bleibt grundsätzlich erhalten (nur im Kooperationsbereich freiwillig eingeschränkt) & geht verloren (einheitliche Leitung bzw. ein einziges Unternehmen)\\
\addlinespace[3pt]
Bindungsintensität & locker bis vertraglich, meist befristet/kündbar & dauerhaft, kapital- oder vertragsbasiert\\
\addlinespace[3pt]
Typische Formen & Arbeitsgemeinschaft/Konsortium, Interessengemeinschaft, Joint Venture, strategische Allianz, Verband, Kartell & Konzern, Fusion (Verschmelzung), Trust\\
\addlinespace[3pt]
Typische Ziele & Synergien nutzen, Kosten teilen, Projekte stemmen --- bei voller Eigenständigkeit & Marktmacht, einheitliche Leitung, volle Kontrolle, Rationalisierung\\
\addlinespace[3pt]
Kartellrechtliche Sicht & grundsätzlich erlaubt --- \emph{außer} wettbewerbsbeschränkende Absprachen (Kartellverbot) & ab bestimmten Größen anmeldepflichtig: Fusionskontrolle (Teil~2)\\
\bottomrule
\end{tabular}
\end{center}

\subsection{Die Systematik als Schaubild}

\begin{center}
\begin{tikzpicture}[
  pfeil/.style={roughline, shorten <=2pt, shorten >=2pt}
]
  \node[roughaccent, font=\footnotesize\bfseries, text=bfwAccent!55!black,
        text width=5.6cm, align=center, minimum height=0.85cm] (root) at (0,2.9)
    {Unternehmenszusammenschlüsse};
  \node[roughbox, text width=6.4cm, align=center, font=\scriptsize] (koop) at (-3.95,0)
    {{\footnotesize\bfseries Kooperation}\\[1pt]
     rechtlich \textbf{und} wirtschaftlich selbstständig\\[4pt]
     Arbeitsgemeinschaft / Konsortium\\
     Interessengemeinschaft\\
     Joint Venture \textperiodcentered{} Strategische Allianz\\
     Verband\\[2pt]
     \textcolor{bfwWarn!85!black}{\bfseries Kartell --- grundsätzlich verboten (\S\,1 GWB)}};
  \node[roughbox, fill=bfwBgAmber, draw=bfwAccent, text width=6.4cm, align=center,
        font=\scriptsize] (konz) at (3.95,0)
    {{\footnotesize\bfseries Konzentration}\\[1pt]
     wirtschaftliche Selbstständigkeit geht verloren\\[4pt]
     Konzern (rechtlich noch selbstständig)\\
     Fusion/Verschmelzung (mind.\ ein Partner erlischt)\\
     Trust (vollständige Verschmelzung)\\[2pt]
     $\rightarrow$ ab Schwellenwerten: Fusionskontrolle (Teil~2)};
  \draw[pfeil] (root.south) -- (koop.north);
  \draw[pfeil] (root.south) -- (konz.north);
\end{tikzpicture}
\end{center}

\begin{merksatz}[title={\bfseries Merksatz: Die Kernabgrenzung}]
\textbf{Kooperation} = freiwillige Zusammen\emph{arbeit} rechtlich und wirtschaftlich
selbstständiger Unternehmen. \textbf{Konzentration} = Zusammen\emph{fassung} von
Unternehmen unter Verlust der wirtschaftlichen Selbstständigkeit --- beim Konzern bleibt
die rechtliche Selbstständigkeit bestehen, bei Fusion und Trust geht auch sie verloren.
\end{merksatz}

\section{Formen der Kooperation}

Die Kooperationsformen unterscheiden sich in ihrer \emph{Bindungsintensität} --- von der
lockeren Interessenvertretung bis zur gemeinsamen Tochtergesellschaft. Jenseits der
gestrichelten Grenze beginnt die Konzentration (Teil~2):

\begin{center}
\begin{tikzpicture}[
  stufe/.style={roughbox, font=\scriptsize\bfseries, text=bfwPrimaryDark,
                text width=1.7cm, align=center, minimum height=1.0cm},
  kstufe/.style={stufe, fill=bfwBgAmber, draw=bfwAccent, text=bfwAccent!55!black}
]
  \node[stufe] (s1) at (0,0)     {Verband};
  \node[stufe] (s2) at (2.2,0.5) {Interessen-\\ gemeinschaft};
  \node[stufe] (s3) at (4.4,1.0) {ARGE /\\ Konsortium};
  \node[stufe] (s4) at (6.6,1.5) {Strategische Allianz};
  \node[stufe] (s5) at (8.8,2.0) {Joint\\ Venture};
  \node[kstufe] (s6) at (11.4,2.5) {Konzern};
  \node[kstufe] (s7) at (13.6,3.0) {Fusion /\\ Trust};
  \draw[roughline, draw=bfwPrimary] (-1.0,-1.0) -- node[below, font=\scriptsize,
    text=bfwInkSoft, pos=0.5]{zunehmende Bindungsintensität $\longrightarrow$} (14.6,-1.0);
  \draw[dashed, thick, draw=bfwWarn] (10.1,-0.9) -- (10.1,3.8)
    node[above, font=\scriptsize, text=bfwWarn!85!black, align=center, pos=1.0]
    {Grenze: Verlust der wirtschaftl.\\ Selbstständigkeit (Teil~2)};
\end{tikzpicture}
\end{center}

\subsection{Arbeitsgemeinschaft (ARGE) und Konsortium}

Rechtlich und wirtschaftlich selbstständige Unternehmen schließen sich \emph{befristet}
zusammen, um \textbf{einen einzelnen Auftrag oder ein Projekt} gemeinsam durchzuführen.
Rechtsform ist regelmäßig die \textbf{GbR}; nach Projektende wird der Zusammenschluss
aufgelöst.

\begin{itemize}
  \item \textbf{Arbeitsgemeinschaft:} typisch in der \emph{Bauwirtschaft} --- mehrere
    Bauunternehmen errichten gemeinsam einen Tunnel oder eine Brücke, weil Kapazität,
    Know-how und Risiko für ein einzelnes Unternehmen zu groß wären.
  \item \textbf{Konsortium:} derselbe Grundgedanke im \emph{Bank- und Finanzwesen} ---
    z.\,B. ein Bankenkonsortium zur gemeinsamen Emission einer Anleihe oder zur
    Finanzierung eines Großkredits.
\end{itemize}

\subsection{Interessengemeinschaft (IG)}

Vertragliche, auf \emph{Dauer} angelegte Zusammenarbeit rechtlich selbstständiger
Unternehmen in \textbf{einzelnen Teilbereichen} --- z.\,B. gemeinsame Forschung und
Entwicklung, gemeinsame Patentverwertung oder ein Gemeinschaftslabor. Häufig wird ein
Gewinn- oder Kostenpool für den gemeinsamen Bereich vereinbart. Die Unternehmen bleiben
im Übrigen unabhängig und treten am Markt weiter getrennt auf.

\subsection{Joint Venture (Gemeinschaftsunternehmen)}

Zwei (oder mehr) Partnerunternehmen gründen ein \textbf{rechtlich selbstständiges
Gemeinschaftsunternehmen}, an dem sie sich mit Kapital beteiligen und dessen Führung und
Risiko sie gemeinsam tragen (oft 50:50). Typische Motive: Erschließung von
\emph{Auslandsmärkten} (Partner bringt Marktkenntnis und Vertriebsnetz ein), Teilung hoher
Entwicklungskosten, Kombination komplementärer Stärken.

\begin{center}
\begin{tikzpicture}[
  partner/.style={roughbox, font=\footnotesize\bfseries, text=bfwPrimaryDark,
                  text width=3.2cm, align=center, minimum height=1.0cm},
  lab/.style={font=\scriptsize, text=bfwInkSoft, fill=white, inner sep=1.5pt}
]
  \node[partner] (a) at (-3.8,2.4) {Partner A\\ \scriptsize\mdseries bleibt selbstständig};
  \node[partner] (b) at (3.8,2.4)  {Partner B\\ \scriptsize\mdseries bleibt selbstständig};
  \node[roughaccent, font=\footnotesize\bfseries, text=bfwAccent!55!black,
        text width=5.4cm, align=center, minimum height=1.1cm] (jv) at (0,0)
    {Joint Venture\\ \scriptsize\mdseries neue, rechtlich selbstständige\\ \scriptsize\mdseries Gesellschaft (z.\,B. GmbH)};
  \draw[roughline, dashed] (a) -- node[lab, above]{Kooperationsvertrag} (b);
  \draw[roughline] (a.south) -- node[lab, below, sloped, align=center, pos=0.45]
    {50\,\% Kapital,\\ Technologie} (jv.west);
  \draw[roughline] (b.south) -- node[lab, below, sloped, align=center, pos=0.45]
    {50\,\% Kapital,\\ Marktzugang} (jv.east);
  \node[font=\scriptsize, text=bfwInkSoft, align=center] at (0,-1.5)
    {Führung und Risiko werden gemeinsam getragen --- typisch für den Eintritt in Auslandsmärkte};
\end{tikzpicture}
\end{center}

\begin{beispiel}
Ein deutscher Automobilhersteller und ein chinesischer Partner gründen gemeinsam eine
Produktions- und Vertriebsgesellschaft in China: Beide Mütter bleiben selbstständig ---
das Joint Venture selbst ist eine neue, eigene Gesellschaft.
\end{beispiel}

\subsection{Strategische Allianz}

Zusammenarbeit --- häufig zwischen (potenziellen) \emph{Konkurrenten} --- in
\textbf{einzelnen strategischen Geschäftsfeldern}, ohne gemeinsames
Tochterunternehmen: z.\,B. Luftfahrt-Allianzen (gemeinsames Streckennetz, Codesharing,
Vielfliegerprogramme) oder Entwicklungs-Partnerschaften bei Batteriezellen. Ziel ist es,
Stärken zu bündeln und Wettbewerbsvorteile gegenüber Dritten zu erlangen; außerhalb der
Allianz bleiben die Partner Wettbewerber.

\subsection{Verband}

Freiwilliger Zusammenschluss zur \textbf{gemeinsamen Interessenvertretung nach außen} ---
gegenüber Staat, Öffentlichkeit, Tarifpartnern --- sowie zur Beratung und Information der
Mitglieder. Beispiele: Branchenverbände (z.\,B. Automobilindustrie), Arbeitgeberverbände.
(Abgrenzung: Die Industrie- und Handelskammern sind \emph{keine} freiwilligen Verbände,
sondern Körperschaften des öffentlichen Rechts mit gesetzlicher Pflichtmitgliedschaft.)
Verbände koordinieren \emph{nicht} das
Marktverhalten ihrer Mitglieder --- sonst drohte ein Kartellverstoß.

\subsection{Kartell --- Kooperation mit kartellrechtlicher Sprengkraft}

Ein \textbf{Kartell} ist eine \emph{vertragliche Absprache} (oder abgestimmte
Verhaltensweise) rechtlich selbstständiger Unternehmen derselben Wirtschaftsstufe mit dem
Ziel, den \textbf{Wettbewerb zu beschränken} --- z.\,B.:

\begin{itemize}
  \item \textbf{Preiskartell:} einheitliche Preise oder Preisuntergrenzen,
  \item \textbf{Quoten-/Produktionskartell:} Aufteilung von Produktionsmengen,
  \item \textbf{Gebietskartell:} Aufteilung von Absatzgebieten oder Kunden,
  \item \textbf{Submissionsabsprachen:} Absprachen bei Ausschreibungen.
\end{itemize}

\begin{beispiel}[title={\bfseries Kartelle in der Praxis --- bekannte Fälle}]
\textbf{Bierkartell (2014):} Das Bundeskartellamt verhängte gegen mehrere große
Brauereien Bußgelder von insgesamt rund 340~Mio.~€ wegen abgestimmter
Fassbier- und Flaschenbierpreise --- ein klassisches \emph{Preiskartell}.
\textbf{LKW-Kartell (2016/17):} Die EU-Kommission belegte führende
LKW-Hersteller mit Geldbußen von zusammen rund 3,8~Mrd.~€ wegen jahrelang
abgestimmter Bruttolistenpreise. \textbf{Kronzeugenregelung:} Wer ein Kartell als
Erster bei der Behörde aufdeckt, kann vollständig bußgeldfrei ausgehen --- so
zerbrechen viele Kartelle von innen.
\end{beispiel}

\begin{achtung}[title={\bfseries Abgrenzung: erlaubte Kooperation oder verbotenes Kartell?}]
Entscheidend ist, ob die Zusammenarbeit eine \emph{Wettbewerbsbeschränkung} bezweckt
oder bewirkt. \textbf{Erlaubt:} ARGE/Konsortium für ein Projekt, das keiner allein
stemmen könnte; gemeinsame Forschung; Verbandsarbeit ohne Marktabsprachen;
freigestellte Mittelstandskooperationen (\S\,2 GWB). \textbf{Verboten:} jede Absprache
über \emph{Preise, Mengen, Gebiete oder Kunden} --- auch informell am Stammtisch oder
per Chat; schon der \emph{Informationsaustausch} über künftige Preise kann genügen.
Prüfungs-Faustregel: Sobald die Absprache den Wettbewerb zwischen den Partnern
ausschaltet, liegt ein Kartell vor.
\end{achtung}

\begin{merksatz}[title={\bfseries Kartellrechtliche Einordnung (Verweis: VWL-Lektion 2)}]
Nach \textbf{\S\,1 GWB} (Gesetz gegen Wettbewerbsbeschränkungen) und
\textbf{Art.\,101 AEUV} sind wettbewerbsbeschränkende Vereinbarungen
\textbf{grundsätzlich verboten} (Verbotsprinzip). Ausnahmen (Freistellung, \S\,2 GWB)
gelten nur, wenn die Vorteile überwiegen --- etwa bei bestimmten
\emph{Mittelstandskooperationen}, die die Wettbewerbsfähigkeit kleiner und mittlerer
Unternehmen verbessern und den Wettbewerb nicht wesentlich beeinträchtigen. Das
\textbf{Bundeskartellamt} verfolgt Verstöße mit empfindlichen Bußgeldern
(Kronzeugenregelung für Aufdecker). Arbeitsgemeinschaften und Konsortien sind dagegen
\emph{grundsätzlich erlaubt}, weil sie den Wettbewerb regelmäßig nicht spürbar
beschränken.
\end{merksatz}

\subsection{Kooperationsformen im Vergleich (Prüfungstabelle)}

\begin{center}
\small
\begin{tabular}{p{3.0cm} p{2.0cm} p{3.3cm} p{2.2cm} p{3.4cm}}
\toprule
\rowcolor{bfwTabHead}\textbf{Form} & \textbf{Dauer} & \textbf{Gegenstand} & \textbf{Rechtsform} & \textbf{Typisches Beispiel}\\
\midrule
ARGE / Konsortium & befristet (ein Projekt) & gemeinsame Auftragsdurchführung & regelmäßig GbR & Tunnelbau-ARGE; Emissionskonsortium der Banken\\
\addlinespace
\hspace{0pt}Interessengemeinschaft & auf Dauer & einzelne Teilbereiche (z.\,B. F\&E, Patente) & Vertrag, ggf.\ Pool & Gemeinschaftslabor zweier Hersteller\\
\addlinespace
Joint Venture & auf Dauer & gemeinsames Tochterunternehmen & neue GmbH/AG & deutsch-chinesisches Produktions-JV\\
\addlinespace
Strategische Allianz & mittel- bis langfristig & einzelne strategische Geschäftsfelder & Vertrag, ohne Tochter & Luftfahrt-Allianz (Codesharing)\\
\addlinespace
Verband & auf Dauer & Interessenvertretung nach außen & e.\,V. & Branchen-, Arbeitgeberverband\\
\addlinespace
Kartell & auf Dauer & \hspace{0pt}Wettbewerbsbeschränkung (Preise, Mengen, Gebiete) & Absprache & \textcolor{bfwWarn!85!black}{verboten, \S\,1 GWB}\\
\bottomrule
\end{tabular}
\end{center}

\section{Formel \& Rechenbeispiel: der Synergieeffekt}

Zum Ziel „Synergien`` gehört ein einfacher Rechenweg. Prägen Sie sich Formel,
Rechenweg und die prüfungsrelevante Schlussfolgerung ein.

\[
\textrm{Synergieeffekt} \;=\; \textrm{Kosten A} + \textrm{Kosten B} - \textrm{Kosten nach der Fusion}
\]

\begin{beispiel}[title={Rechenbeispiel: Kostensynergie}]
A hat Jahreskosten von 60~Mio.~€, B von 40~Mio.~€. Nach der Fusion betragen die
Gesamtkosten dank gemeinsamer Verwaltung und gemeinsamen Einkaufs nur noch 90~Mio.~€.
\begin{enumerate}
  \item \textbf{Kosten vor der Fusion:} $60 + 40 = 100$~Mio.~€
  \item \textbf{Kosten nach der Fusion:} $90$~Mio.~€
  \item \textbf{Synergieeffekt:} $100 - 90 = 10$~Mio.~€ Ersparnis
  \item \textbf{In Prozent:} $10 \div 100 \times 100\,\% = 10\,\%$ der Ausgangskosten
\end{enumerate}
\textbf{Schlussfolgerung:} Das ist das Prinzip „2 + 2 = 5`` --- gemeinsam entsteht mehr
Wert (bzw. weniger Kosten) als die Summe der Einzelteile.
\end{beispiel}

\begin{merksatz}[title={\bfseries Wichtig für die Prüfungsvorbereitung (Teil 1)}]
Sie sollten sicher können: die Abgrenzung \emph{Kooperation vs. Konzentration} über die
rechtliche und wirtschaftliche Selbstständigkeit, die \emph{Kooperationsformen mit
Beispiel} (ARGE/Konsortium, IG, Joint Venture, strategische Allianz, Verband), das
\emph{Kartellverbot} (\S\,1 GWB) mit den Kartellarten und Ausnahmen (\S\,2 GWB) sowie den
Rechenweg \emph{Synergieeffekt}. Weiter geht es in Teil~2 mit Konzern, Fusion, Trust,
den Richtungen und der Fusionskontrolle --- dort auch die Rechenwege Marktanteil und
Beteiligungsquote.
\end{merksatz}

\vspace{1em}
\hrule
\vspace{0.8em}
\noindent\textsf{\small Quellen: Rahmenplan Wirtschaftsbezogene Qualifikationen (DIHK), Abschnitt 1.4;
Übungsband Betriebswirtschaft (DIHK-Bildungs-GmbH), Übungen 54--58; GWB (Grundzüge).
Textband-Auszug wird nachgereicht. Fortsetzung in Teil~2 (Session 2, 10:55--13:15 Uhr).}

\end{document}
